\section{Training Dynamics and Serving Evidence}
\label{sec:systems}
\subsection{RQ3: optimization progress versus useful quality}
Figure~\ref{fig:loss} shows all available loss records for four runs. Shared v1 and extraction use trainer-state records; functional and single-token relevance use the rounded values printed in console logs at ten-step intervals. The functional loss approaches zero while test behavior remains imperfect and scenario-dependent. Extraction loss also falls rapidly, whereas its monitored color accuracy varies between checkpoints. These observations motivate evaluating task behavior throughout training rather than using decreasing loss as a proxy for generalization.

\begin{figure}[t]
\centering
\begin{tikzpicture}
\begin{groupplot}[group style={group size=2 by 2,horizontal sep=1.3cm,vertical sep=1.5cm},width=.46\textwidth,height=4cm,xlabel={Optimizer step},ylabel={Training loss},ymode=log,grid=major,tick label style={font=\scriptsize},label style={font=\small},title style={font=\small}]
\nextgroupplot[title={Shared v1 (trainer state)}]
\addplot[blue!70!black,thick] table[x=step,y=loss,col sep=comma]{evidence/loss_shared_v1.csv};
\nextgroupplot[title={Extraction (trainer state)}]
\addplot[teal!70!black,thick] table[x=step,y=loss,col sep=comma]{evidence/loss_understand.csv};
\nextgroupplot[title={Functional (rounded console log)}]
\addplot[orange!85!black,thick] table[x=step,y=loss,col sep=comma]{evidence/loss_functional.csv};
\nextgroupplot[title={Relevance letters (rounded log)}]
\addplot[purple!70!black,thick] table[x=step,y=loss,col sep=comma]{evidence/loss_single_token.csv};
\end{groupplot}
\end{tikzpicture}
\caption{All saved ten-step training-loss records for the four runs with available loss traces. Each panel has its own logarithmic vertical scale. Loss magnitudes differ across objectives and must not be ranked as task quality. Decreasing functional loss is compatible with limited generalization and repeated exposure to 131 training rows.}
\label{fig:loss}
\end{figure}


Loss magnitudes are not directly comparable across panels: objectives differ in token weighting, completion length, and data distribution. The functional log reports approximately 389.2 seconds of trainer runtime and the single-token relevance log approximately 1,004 seconds. These are measurements from individual runs, not end-to-end campaign cost, and do not include every failed experiment, provisioning interval, upload, API call, or idle period. Exact GPU billing and a complete per-run environment ledger are absent. We therefore do not derive a dollars-per-quality or hardware Pareto frontier from these durations.

\subsection{Implemented service and observed load-test limits}
The FastAPI implementation has relevance, identity, functional-relation, and catalog-normalization routes. It lazily loads three separate base-plus-adapter model instances for shared Match, functional relation, and Understand. The process has per-loader locks intended to prevent duplicate concurrent loading. This is different from keeping one resident base and switching adapters, and no memory comparison between the two strategies was measured. A technical-compatibility inference route is not present in the inspected API.

The historical retained Locust CSVs are decisive about what was measured (Table~\ref{tab:serving}). Before the loading-lock change, the statistics contain one completed health request and no completed model inference. The after-change CSV contains zero completed requests. The accompanying narrative describes a three-user, 90-second CPU test and an isolated slow functional response, but the raw server traces needed to reproduce the loading diagnosis are not archived. The code contains the lock; the CSVs do not demonstrate recovered inference throughput.

\begin{table}[t]
\centering\small
\caption{Completed requests in retained load-test statistics. No model-latency percentile or inference throughput is estimable from these completions. Zero logged failures does not imply that pending requests succeeded.}
\label{tab:serving}
\begin{tabular}{lrrrl}
\toprule
Artifact & Health completed & Inference completed & Logged failures & Inference p50/p99 \\
\midrule
Before lock & 1 & 0 & 0 & Not observed \\
After lock & 0 & 0 & 0 & Not observed \\
\bottomrule
\end{tabular}
\end{table}

A health response checks adapter-path availability, not a completed warm inference or a readiness guarantee. Requests that remain unfinished at the end of a test are censored observations; assigning them zero latency, zero error probability, or the health route's latency would be incorrect. No GPU concurrency curve, sustained request-rate measurement, completed-request tail latency, or production availability estimate is established by these artifacts.

\subsection{Cost accounting must include successful task completion}
The service exposes an estimate based on a configured \$0.74/hour rate and an assumed 20 tokens/second. Its resulting \$0.0103 per 1,000 tokens is arithmetic over assumptions, not an observed serving cost. Training iterations per second do not validate decoding tokens per second, and prompt processing and generation cannot generally be priced by a single assumed throughput. We exclude this estimate from comparative results.

For a future serving evaluation, a useful task-level quantity is
\begin{equation}
 C_{\mathrm{valid,SLO}}=
 \frac{C_{\mathrm{GPU}}+C_{\mathrm{CPU}}+C_{\mathrm{storage}}+C_{\mathrm{network}}+C_{\mathrm{API}}}
 {\sum_{i=1}^{N}\mathbb{1}[\text{request }i\text{ correct, schema-valid, and within SLO}]},
\label{eq:cost}
\end{equation}
where all billed resources, including loading and idle capacity, are counted over the same declared interval. A zero denominator is undefined rather than a favorable estimate. Reporting this alongside raw cost/request, completion rate, p50/p95/p99 latency, and throughput separates low unit price from useful service. Equation~\ref{eq:cost} is a proposed accounting criterion; the historical study does not supply the measurements needed to instantiate it. The new HTTP prototype reports a narrower GPU-rate estimate, not a complete invoiced SLO cost.

\section{Discussion, Limitations, and Responsible Use}
\label{sec:discussion}
\paragraph{What the evidence supports.}
Small-model specialization appears promising for this narrow extraction sample: it produces valid JSON and favorable field F1 relative to the measured off-the-shelf and API baselines. The matching experiments demonstrate that one shared adapter can emit several task labels, but they do not establish a uniformly strong horizontal commerce model. The letter-label relevance analysis illustrates why measuring all classes is necessary even after obvious output collapse disappears. The functional analysis shows that good development accuracy can coexist with scenario reuse and test-informed selection.

\paragraph{What makes the analysis useful beyond this artifact.}
Four failures arise at different levels: shared wording creates scorer collisions; paraphrases share scenario-level information; post-fit repartitioning fails to remove prior exposure; and aggregate score tables lose the pairing required for inference. Their remedies require different units of control. A strict parser cannot fix scenario overlap, and more data cannot repair an exposed test set. The comparison contract makes these dependencies explicit before a result is promoted into a model-selection or deployment claim. The aggregate-pairing calculation is a compact diagnostic of an additional information loss: even correct marginal counts can leave a significance conclusion unidentified.

\paragraph{Limits to generalization and novelty.}
The study has one base size, no repeated seeds, limited category coverage, and no measured multilingual, merchant-held-out, temporal, or adversarial robustness. Unknown/abstention behavior and calibration are not evaluated. The underlying public corpora may have appeared during pretraining; that exposure is unmeasured. The strongest comparison uses a permissive scorer and a small partially overlapping sample. The synthetic judge is neither blinded nor human, and no naturalness annotation is available. No claim of being the first commerce SLM, of beating all related Hugging Face models, or of inventing multitask QLoRA follows from this evidence. Novel training, data-selection, and serving mechanisms would require implementation and controlled evaluation beyond the present record.

\paragraph{Experiments needed to strengthen the result.}
The highest-priority extension is a new, frozen evaluation suite with query-, entity-, and scenario-level separation as appropriate; strict label parsers; exact and normalized extraction metrics; and retained per-example outputs. Shared and specialist models should then be trained with matched data, budgets, and multiple seeds, alongside an unfine-tuned Qwen baseline and supervised compact alternatives. A second phase should vary real/synthetic mixture, scenario grouping, label representation, and task balancing one factor at a time. Finally, GPU serving measurements should evaluate quality-preserving precision and batching choices under declared service-level objectives. The new development and serving follow-up covers part of this agenda; controlled retraining, human review, full comparator coverage and final evaluation remain prospective. Appendix~\ref{app:future} specifies the comparisons in more detail.

\paragraph{Responsible use and data rights.}
Product matching errors can merge distinct offers or misrepresent fit; functional complementarity must not be treated as a guarantee of technical compatibility. The present diagnostics do not justify autonomous compatibility assurance. The input corpora carry their respective publisher terms, and a repository manifest is not a substitute for reviewing redistribution and commercial-use conditions. The experiments use product data and synthetic descriptions rather than customer conversations; that does not establish a comprehensive privacy audit. No human annotation study was conducted in this work. Deployment would require task-specific validation, evidence-aware handling of uncertain outputs, and appropriate operational controls; these capabilities are not established by a request schema containing a tenant identifier.

\paragraph{Research assistance disclosure.}
Generative AI assisted the manuscript revision, code inspection, literature synthesis, and implementation of the offline evidence audit. Synthetic training-data generation and automated judging are described separately in Section~\ref{sec:data}. The retrospective analyses use preserved artifacts; the follow-up retains newly generated predictions and request logs. Responsibility for validating the manuscript and any eventual submission remains with the author.

\section{Conclusion}
\method{} provides a small-model adaptation case study whose strongest supported signal is brand/color extraction on a common 200-listing sample. Reconstructing the full evidence changes the interpretation of the broader campaign: matching comparisons are constrained by unequal samples, minority diagnostics reuse scenarios, checkpoint selection exposes a nominal test partition, and historical serving traces contain no completed inference load measurements. The separately documented follow-up adds matched development comparisons and a bounded GPU serving prototype without converting them into final-test claims. The released audit makes these boundaries reproducible and adds class-level analysis, scorer checks, and aggregate paired-uncertainty bounds. A credible path toward a broadly useful commerce model requires preserving evaluation independence and operational measurements as carefully as the model weights themselves.
